\section*{Capítulo I, Problema VII}

\textbf{Problema:}

7. Considera funciones continuas en el intervalo $[-\pi, \pi]$. Define un producto escalar similar al anterior para este intervalo. Demuestra que las funciones $\sin(nx)$ y $\cos(mx)$ son ortogonales para este producto escalar ($m, n$ siendo enteros).

\textbf{Solución:}

El producto escalar para funciones continuas en el intervalo $[-\pi, \pi]$ está definido como:
\[
    \langle f, g \rangle = \int_{-\pi}^{\pi} f(x)g(x) \, dx.
\]

Para demostrar que $\sin(nx)$ y $\cos(mx)$ son ortogonales, calculamos su producto escalar:
\[
    \langle \sin(nx), \cos(mx) \rangle = \int_{-\pi}^{\pi} \sin(nx) \cos(mx) \, dx.
\]

Usamos la identidad trigonométrica:
\[
    \sin(nx)\cos(mx) = \frac{1}{2}\left[\sin((n+m)x) + \sin((n-m)x)\right].
\]

Sustituyendo en el producto escalar:
\[
    \langle \sin(nx), \cos(mx) \rangle = \frac{1}{2} \int_{-\pi}^{\pi} \sin((n+m)x) \, dx + \frac{1}{2} \int_{-\pi}^{\pi} \sin((n-m)x) \, dx.
\]

Sabemos que la integral de $\sin(kx)$ sobre un intervalo simétrico $[-a, a]$ es cero para cualquier entero $k \neq 0$. Por lo tanto:
\[
    \langle \sin(nx), \cos(mx) \rangle = 0.
\]

Esto demuestra que $\sin(nx)$ y $\cos(mx)$ son ortogonales bajo el producto escalar definido en el intervalo $[-\pi, \pi]$. 
